\documentclass[11pt,a4paper]{article}
\usepackage[T1]{fontenc}
\usepackage[utf8]{inputenc}
\usepackage{lmodern}
\usepackage[margin=2.6cm]{geometry}
\usepackage{amsmath,amssymb,amsthm}
\usepackage{booktabs}
\usepackage{array}
\usepackage{enumitem}
\usepackage[hidelinks]{hyperref}

\newtheorem{theorem}{Theorem}
\newtheorem{proposition}[theorem]{Proposition}
\newtheorem{lemma}[theorem]{Lemma}
\newtheorem*{corollary}{Corollary}
\newtheorem*{lemmaS}{Lemma S}
\newtheorem*{lemmaP}{Lemma P}
\newtheorem*{avis}{Theorem}
\newtheorem*{fact}{Fact (from C1, C2)}
\theoremstyle{definition}
\newtheorem*{definition}{Definition}
\newtheorem{question}{Question}
\newtheorem*{questionz}{Question}
\theoremstyle{remark}
\newtheorem*{remark}{Remark}
\newtheorem*{remarks}{Remarks}

\newcommand{\MET}{\mathrm{MET}}
\newcommand{\cl}{\mathrm{cl}}
\newcommand{\Z}{\mathbb{Z}}
\newcommand{\Q}{\mathbb{Q}}
\newcommand{\R}{\mathbb{R}}
\newcommand{\eps}{\varepsilon}
\newcommand{\file}[1]{\texttt{#1}}
\newcommand{\comp}[1]{\medskip\noindent\textbf{#1}}
\newcommand{\step}[1]{\smallskip\noindent\emph{#1}}

\title{No bounded-scope characterisation of the extreme rays of the metric cone}
\author{Adrian Burlacu}
\date{Draft of 26 November 2026}

\begin{document}
\maketitle

\begin{abstract}
A metric on $n$ points spans an extreme ray of the metric cone $\MET_n$ exactly when its tight triangle equalities
have rank $\binom n2 - 1$. We ask whether this rank condition can be replaced by a test that looks at a bounded
number of points at a time. To measure this we introduce the \emph{locality} $L(d)$ of an extreme ray $d$: the least
$s$ such that extremality of $d$ can be derived by solving the tight equalities inside sets of $s$ points and
propagating the proportionalities obtained. We prove that $L = 4$, the smallest possible value, for every graph
covered by Avis's isometric-cycle criterion, hence for almost all graphs, and that $L$ is nevertheless unbounded: an
explicit chiral family $C_t$ of extreme rays on $2t$ points ($t$ odd, $t \ge 9$) has $L(C_t) = 2t$, that is, no
proper subset of the points proves extremality (the bound $L(C_t) > t/6$ has a short proof; the equality is
computer-assisted, with exact arithmetic). Hence no criterion of bounded scope characterises the extreme rays of the
metric cone, or the vertices of the metric polytope $m_n$ adjacent to a cut. We compute $L$ exactly for all extreme
rays of $\MET_n$, $n \le 8$, and for the 1,055,145 cut-adjacent orbits of the known vertex list of $m_9$: the largest
values are $4, 4, 5, 7$ for $n = 5, \dots, 8$ and $9$ on the $m_9$ list, and the share of rays with $L = 4$ falls
from 97\% ($n = 7$) to 30\% ($n = 9$). We also show, by a counting argument resting on a theorem of Avis, that a
generator of vertices of $m_n$ which makes fewer than $(1 - \delta)n$ bits of choice when it adds the $n$-th node
produces a vanishing share of the orbits, so that no finite set of such rules with a finite list of exceptions
describes the vertices.
\end{abstract}

Every statement below is marked \textbf{Theorem} / \textbf{Proposition} / \textbf{Lemma} (proved here),
\textbf{Computation} (items C1--C12: exact computer results on a stated data set) or \textbf{Question} (open).

\section{Introduction}

The metric cone $\MET_n$ is the set of $x \in \R^{\binom n2}$ with $x_{ij} \le x_{ik} + x_{jk}$ for all triples, and
the metric polytope $m_n$ is its intersection with the half-spaces $x_{ij} + x_{ik} + x_{jk} \le 2$. The cut polytope
is contained in $m_n$ and has the same integral vertices, and $m_n$ is the standard linear relaxation of max-cut; its
fractional vertices are what separates the relaxation from the problem. Complete vertex lists are known for
$n \le 8$ \cite{DFPS01,DFMV03}, and for $n = 9$ there is a list of 1,056,368 orbits conjectured to be complete
\cite{DI07}. Constructions of infinite families are known (cuts and anticuts, trivial extensions, graphic vertices,
Avis's doubling, gluing; see \cite{LP92,L96,DL97}), but they account for a small and shrinking part of the lists, and
no description of the vertices other than enumeration is known.

\begin{questionz}
Is there a finite set of rules, valid for every $n$, and a finite list of exceptions, which together produce exactly
the vertices of $m_n$?
\end{questionz}

The answer depends on what a rule may do, and we separate two regimes.

\emph{Graph-input rules.} A rule receives a weighted graph on $n$ nodes and returns its shortest-path metric, scaled.
Every vertex adjacent to a cut arises this way, and the rule is correct exactly when the metric is an extreme ray of
$\MET_n$, that is, when its tight triangle equalities have rank $\binom n2 - 1$. The question becomes whether this
rank condition can be replaced by a combinatorial test of bounded scope, in the way Avis's isometric-cycle criterion
\cite{A80} does for a large class of graphs. We measure the scope a test needs by the locality $L$
(Section~\ref{sec:locality}).

\emph{Small-choice rules.} A rule receives an object on $n - 1$ nodes and a bounded or slowly growing amount of
choice and returns an object on $n$ nodes. All classical constructions are of this kind.

\paragraph{Results.}
\begin{enumerate}[leftmargin=*]
\item \emph{Locality.} The invariant $L$ is well defined, sound and complete (Proposition~\ref{prop:closure}). It
  equals 4 for every graph covered by the isometric-cycle criterion, and for almost all graphs
  (Theorem~\ref{thm:four}).
\item \emph{Locality is unbounded} (Theorem~\ref{thm:main}, the main result). An explicit family $C_t$ of extreme
  rays on $n = 2t$ points has $L = n$. The proof reduces the tight system of $C_t$ to a linear recurrence around the
  cycle whose characteristic polynomial has a root $\lambda$ that is not a root of unity. On any set of points that
  misses five consecutive positions this root yields a second solution of the local equations, which no local
  deduction can exclude; this gives $L > n/12$ by hand. When a single point is removed, the recurrence is broken at
  that point only, and a solution that decays like $\lambda^{-k}$ away from the break still exists; showing that it
  separates all classes is a finite exact computation, uniform in $t$.
\item \emph{Counting.} The number of orbits of vertices of $m_n$ is at least $2^{n^2/2 - O(n \log n)}$
  (Proposition~\ref{prop:count}, from Avis's density theorem), and a generator with at most $(1 - \delta)n$ bits of
  choice per node produces a vanishing share of them (Theorem~\ref{thm:generator}).
\item \emph{Exact computations.} $L$ for all extreme rays of $\MET_n$, $n \le 8$, and for the $m_9$ list (C1--C3);
  searches for large locality beyond the lists (C4--C6); and negative results for four other candidate descriptions
  of the $m_9$ list (C7--C12).
\end{enumerate}

Together, 1 and 2 say that the scope of the rank condition is the smallest possible for a typical graph and
unbounded in the worst case; the computations show that on the actual vertex lists the typical case is already
leaving the minimum at $n = 9$.

\section{Preliminaries}

\paragraph{Symmetries.} The symmetry group of $m_n$ is generated by the permutations of the nodes and the switchings
$r_S(x)_{ij} = 1 - x_{ij}$ for $ij$ in the cut $\delta(S)$, $x_{ij}$ otherwise. It has order $n!\cdot 2^{n-1}$ for
$n \ge 5$. ``Orbit'' means an orbit of this group. The cuts are the integral vertices and form one orbit.

\paragraph{Tight systems.} For a metric $d$ on $n$ points, a triple $(a, m, b)$ is \emph{tight} if
$d_{ab} = d_{am} + d_{mb}$ ($m$ is between $a$ and $b$). The \emph{tight system} $T(d)$ is the set of linear
equations $x_{ab} = x_{am} + x_{mb}$ over the tight triples. $d$ spans an extreme ray of $\MET_n$ iff the solution
space of $T(d)$ is the line through $d$, i.e.\ $\operatorname{rank} T(d) = \binom n2 - 1$.

\begin{lemma}[rays give vertices]\label{lem:vertex}
Let $d$ span an extreme ray of $\MET_n$, $n \ge 3$, and let $\tau$ be the largest perimeter
$d_{ij} + d_{ik} + d_{jk}$. Then $x = (2/\tau)\cdot d$ is a vertex of $m_n$, and it is adjacent to the cut vertex $0$.
\end{lemma}

\begin{proof}
$x \in m_n$ by the choice of $\tau$. The equations of $T(d)$ are tight at $x$ and have rank $\binom n2 - 1$; each of
them vanishes at $x$. A triple attaining $\tau$ gives a tight perimeter inequality $x_{ij} + x_{ik} + x_{jk} = 2$,
which does not vanish at $x$ and is therefore independent of them. So the tight inequalities at $x$ have rank
$\binom n2$. The segment from $0$ to $x$ lies on the $\binom n2 - 1$ independent facets of $T(d)$, so it is an edge.
\end{proof}

Conversely, a vertex of $m_n$ adjacent to a cut is, after the switching that moves this cut to $0$, of the form
$(2/\tau)d$ for an extreme ray $d$ \cite{LP92}. Every vertex of $m_n$, $n \le 8$, is adjacent to a cut; on the $m_9$
list 1,221 orbits are not \cite{DI07} (re-derived in C7 below).

\paragraph{Weighted graphs.} For a metric $d$ with positive distances, the \emph{minimal graph} has an edge $ab$ of
weight $d_{ab}$ whenever no third point is between $a$ and $b$; $d$ is its shortest-path metric. We scale $d$ to
coprime integers and call the set of weights the \emph{alphabet}. \emph{Graphic} means alphabet $\{1\}$; $d_G$
denotes the shortest-path metric of a connected graph $G$.

\section{Locality of an extreme ray}\label{sec:locality}

Let $d$ be a metric on a node set $V$, $|V| = k \ge 4$, with positive distances. Write a solution $x$ of $T(d)$ as
$x_p = d_p\cdot u_p$ over the pairs $p$; then $u \equiv 1$ is a solution, and $d$ is extreme iff every solution has
$u$ constant.

\begin{definition}[level-$s$ closure, locality]
Let $\pi$ be a partition of the pairs, and $W \subseteq V$. The system $T_W(\pi)$ consists of the equations
$d_{ab}\cdot u_{[ab]} = d_{am}\cdot u_{[am]} + d_{mb}\cdot u_{[mb]}$ for the tight triples with $a, m, b \in W$, in
one variable $u_C$ per class $C$ of $\pi$ that contains a pair inside $W$. $W$ \emph{forces} $C = C'$ under $\pi$ if
$u_C = u_{C'}$ in every solution of $T_W(\pi)$. The \emph{level-$s$ closure} $\cl_s(d)$ is the finest partition
$\pi$ that is closed under merging $C$ and $C'$ whenever some $W$ with $|W| = s$ forces $C = C'$ under $\pi$. The
\emph{locality} $L(d)$ is the least $s \ge 4$ such that $\cl_s(d)$ has a single class; $L(d) = \infty$ if there is
none.
\end{definition}

\begin{proposition}\label{prop:closure}
\begin{enumerate}[label=\textup{(\roman*)},leftmargin=*]
\item $\cl_s(d)$ is well defined, and is the result of applying the merges in any order until none applies.
\item (Soundness) If two pairs $p, q$ are in the same class of $\cl_s(d)$, then $x_p/d_p = x_q/d_q$ for every
  solution $x$ of $T(d)$. Hence $L(d) < \infty$ implies that $d$ is extreme.
\item $\cl_s(d)$ refines $\cl_{s+1}(d)$, and $\cl_k(d)$ has a single class iff $d$ is extreme. Hence $d$ is extreme
  iff $L(d) \le k$, and then $4 \le L(d) \le k$.
\item $L$ is invariant under relabelling and scaling.
\end{enumerate}
\end{proposition}

\begin{proof}
(i) If $\pi$ refines $\pi'$, every solution of $T_W(\pi')$ is a solution of $T_W(\pi)$ that is constant on the
classes of $\pi'$; so a merge forced under $\pi$ is forced, or already made, under $\pi'$. The merge operator is
therefore monotone on the lattice of partitions and has a least closed partition above the discrete one, reached by
any maximal sequence of merges. (ii) By induction on the sequence of merges: if $u$ is constant on the classes of
$\pi$ for every solution of $T(d)$, then the restriction of $u$ to the pairs inside $W$ solves $T_W(\pi)$, so a
forced merge preserves the property. (iii) A set $W$ of size $s$ is contained in a set $W'$ of size $s + 1$ and
$T_W(\pi) \subseteq T_{W'}(\pi)$, which gives the refinement. For $W = V$ and $\pi$ discrete, $T_V(\pi)$ is $T(d)$
in the variables $u$, and it forces all classes equal iff its solution space is spanned by $u \equiv 1$. (iv) is
clear.
\end{proof}

\paragraph{What $L$ bounds.} Call a deduction \emph{$s$-local} if it derives a proportionality $x_p = c\cdot x_q$
from the tight equations among at most $s$ nodes, used together with proportionalities derived earlier. By (i), the
proportionalities derivable by $s$-local deductions are exactly those of $\cl_s(d)$. So a test for extremality that
proceeds by $s$-local deductions succeeds on $d$ iff $L(d) \le s$, and a characterisation of extreme rays by
forbidden or required configurations on at most $s$ nodes, closed under such propagation, cannot be complete unless
$L \le s$ on all rays. The notion covers deductions of two-term proportionalities; a local rule that carries longer
linear relations from one subset to another is not covered and would need a wider definition.

\begin{theorem}\label{thm:four}
Let $G$ be a connected graph on $k \ge 4$ nodes.
\begin{enumerate}[label=\textup{(\alph*)},leftmargin=*]
\item If all edges of $G$ fall into one class under the relation generated by ``opposite edges of an even isometric
  cycle'' ($G$ is 1-ic-colorable in the sense of Avis \cite[\S2]{A80}), then $L(d_G) = 4$; in particular $d_G$ is
  extreme.
\item For $G = G(n, 1/2)$, with probability $1 - o(1)$: $G$ is connected and any two edges of $G$ are linked by a
  chain of induced 4-cycles in which consecutive edges are opposite. Hence $L(d_G) = 4$ for
  $2^{\binom n2}\cdot(1 - o(1))$ labelled graphs on $n$ nodes.
\end{enumerate}
\end{theorem}

\begin{proof}
(a) Let $1, 2, \dots, 2t$ be an isometric cycle of $G$, $t \ge 2$, and $W = \{1, 2, t + 1, t + 2\}$. The triples
$(1, t + 2, t + 1)$, $(1, 2, t + 1)$, $(2, 1, t + 2)$, $(2, t + 1, t + 2)$ (middle node between the outer two) are
tight and lie in $W$. They give $x_{1,t+1} = x_{1,t+2} + x_{t+1,t+2} = x_{1,2} + x_{2,t+1}$ and
$x_{2,t+2} = x_{1,2} + x_{1,t+2} = x_{2,t+1} + x_{t+1,t+2}$; subtracting the two resulting equalities gives
$x_{1,2} = x_{t+1,t+2}$ (this is \cite[Lemma 2.3]{A80}). Both pairs have $d = 1$, so $W$ forces their classes equal
under any partition. Hence all edges lie in one class $E$ of $\cl_4(d_G)$. Let $ab$ be a pair at distance $j \ge 2$
and $m$ a neighbour of $a$ on a geodesic from $a$ to $b$. If $[mb] = E$ (induction on $j$), the equation
$j\cdot u_{ab} = u_E + (j - 1)\cdot u_E$ inside any 4-set $W \supseteq \{a, m, b\}$ forces $[ab] = E$. So $\cl_4$ has
one class, $L \le 4$, and $L \ge 4$ by definition; extremality follows from Proposition~\ref{prop:closure}(ii),
without using \cite[Thm 2.4]{A80}.

(b) Fix distinct nodes $s, t, u, v$ and split the other $n - 4$ nodes into $\lfloor (n - 4)/2 \rfloor$ pairs
$\{x, y\}$. For a pair, let $A_{x,y}$ be the event that $xs, xy, xv, yt, yu$ are edges and $xt, xu, ys, yv$ are not;
it has probability $2^{-9}$, and the events for different pairs are independent of each other and of the presence of
$st$ and $uv$, since they concern disjoint sets of pairs of nodes. If $st$ and $uv$ are edges and $A_{x,y}$ holds,
then $s, t, y, x$ and $u, v, x, y$ are induced 4-cycles, in which $st$ and $uv$ are both opposite to $xy$
\cite[Fig.~3.5]{A80}. There are at most $3\binom n4$ choices of two disjoint pairs $st$, $uv$, so the probability
that some two disjoint edges are not linked in this way is at most
$3\binom n4 (1 - 2^{-9})^{\lfloor (n-4)/2 \rfloor} \to 0$. Two edges with a common endpoint are linked through a
third edge disjoint from both, which exists unless some $n - 3$ nodes span no edge, an event of probability at most
$\binom n3\cdot 2^{-\binom{n-3}{2}} \to 0$. The probability of an isolated node is at most $n\cdot 2^{-(n-1)} \to 0$,
and a graph without isolated nodes in which all edges are linked by chains of 4-cycles is connected. An induced
4-cycle is isometric, so (a) applies.
\end{proof}

Part (b) is the case $p = 1/2$ of Avis's density theorem \cite[Thm 3.3]{A80}, with the stronger conclusion $L = 4$;
the argument follows Avis's and treats explicitly the pairs of edges with a common endpoint.

\begin{remark}[finite $n$]
The union bound in (b) is effective only for $n$ in the thousands, but the conclusion sets in much earlier (C4
below): of 300 samples of $G(n, 1/2)$ each, the shares with $d_G$ extreme are 0.7\%, 6\%, 23\%, 61\%, 81\%, 91\%,
99.3\%, 100\% for $n = 6, 8, 10, 12, 14, 16, 20, 24$, and all but one of these 1,383 rays have $L = 4$.
\end{remark}

\section{Locality is unbounded}\label{sec:main}

\paragraph{The family $C_t$.} For odd $t \ge 9$ let $C_t$ be the following function on the $2t$ points $A_i, B_i$
($i \in \Z_t$), where $|i - j|$ denotes the cyclic distance and offsets $o$ are taken in $(-t/2, t/2)$:
\begin{center}
\begin{tabular}{ll}
\toprule
pair & distance \\
\midrule
$A_i, A_j$ & $16$ if $|i - j| = 1$;\quad $12$ if $|i - j| = 2$;\quad $8$ if $|i - j| \ge 3$ \\
$B_i, B_j$ & $10$ if $|i - j| = 1$;\quad $20$ if $|i - j| = 2$;\quad $22$ if $|i - j| \ge 3$ \\
$A_i, B_{i+o}$ & $13, 11, 7, 9, 19$ for $o = -2, -1, 0, 1, 2$;\quad $15$ if $|o| \ge 3$ \\
\bottomrule
\end{tabular}
\end{center}
$C_t$ is invariant under the rotation $i \mapsto i + 1$ and has no reflection symmetry. $C_\infty$ denotes the same
rule on the line ($i \in \Z$, ordinary offsets).

\begin{theorem}\label{thm:main}
\begin{enumerate}[label=\textup{(\alph*)},leftmargin=*]
\item For every odd $t \ge 9$, $C_t$ is a metric and spans an extreme ray of $\MET_{2t}$.
\item For every odd $t \ge 9$ and every $s \le t/6$, the level-$s$ closure of $C_t$ is not total. Hence
  $L(C_t) > t/6$ and
  \[
    \sup\,\{\, L(d) : d \text{ an extreme ray of } \MET_n \,\} > n/12 \qquad\text{for every } n = 2t,\ t \text{ odd},\ t \ge 9;
  \]
  in particular it is unbounded.
\item $L(C_t) = 2t$ for every odd $t \ge 9$: no proper subset of the points proves extremality. (Computer-assisted:
  for $t \le 21$ by exact computation of the level-$(2t - 1)$ closure; for $t \ge 23$ by the proof below, which uses
  exact finite checks in $\Q(\lambda)$.)
\end{enumerate}
\end{theorem}

\begin{corollary}
No test that proceeds by $s$-local deductions, for any fixed $s$, characterises the extreme rays of the metric cone;
by Lemma~\ref{lem:vertex} the same holds for the vertices of $m_n$ adjacent to a cut. Moreover
$\max\,\{ L(d) : d \text{ extreme in } \MET_n \} = n$ for every $n \equiv 2 \pmod 4$, $n \ge 18$.
\end{corollary}

\begin{proof}[Proof of (a)]
\step{Patterns.} A distance of $C_t$ depends on the kinds of the two points and on their offset capped at $\pm 3$.
For $t \ge 9$ every configuration of three points, described by kinds and capped offsets, occurs already in $C_9$
(two of the three cyclic gaps determine the capped configuration, and all combinations with gaps $0$, $1$, $2$ or
``at least 3'' are realised when $t = 9$), and the same configurations occur in $C_\infty$. Therefore the triangle
inequalities, and the list of types of tight triples, need to be checked only for $t = 9$; the check (V1 in
\file{family\_C\_proof.py}) shows that $C_9$ is a metric and that there are 11 types of tight triples, written
long = short + short:
\begin{align*}
& A_iA_{i+1} = A_iA_k + A_kA_{i+1}\ (k \text{ far from both}); \qquad A_iA_{i+1} = A_iB_{i+1} + B_{i+1}A_{i+1};\\
& A_iB_{i+2} = A_iA_{i+2} + A_{i+2}B_{i+2} = A_iA_{i+3} + A_{i+3}B_{i+2} = A_iB_{i+1} + B_{i+1}B_{i+2};\\
& A_iB_j = A_iA_j + A_jB_j\ (|i - j| \ge 3);\\
& B_iB_{i+2} = B_iA_{i+2} + A_{i+2}B_{i+2} = B_iA_{i+1} + A_{i+1}B_{i+2} = B_iB_{i+1} + B_{i+1}B_{i+2};\\
& B_iB_j = B_iA_i + A_iB_j\ (|i - j| \ge 3;\text{ also with $i$ and $j$ exchanged});\\
& B_iB_{i+3} = B_iA_{i+2} + A_{i+2}B_{i+3}.
\end{align*}

\step{The $A$-pairs.} Let $x$ be a solution of the tight system. Write $f(i, k) = x(A_iA_k)$ for $|i - k| \ge 3$ and
$\rho_i = x(A_iA_{i+1})$. By the first type, $\rho_i = f(i, k) + f(k, i + 1)$ whenever $k - i \in [4, t - 3]$. For
$k - i \in [4, t - 4]$ this gives four equations for the two pairs $\{i, i+1\}$ and $\{k, k+1\}$:
$\rho_i = f(i,k) + f(i+1,k) = f(i,k+1) + f(i+1,k+1)$ and $\rho_k = f(i,k) + f(i,k+1) = f(i+1,k) + f(i+1,k+1)$, hence
$f(i+1, k) = f(i, k+1)$ and $f(i, k) = f(i+1, k+1)$. The second identity shows that $f(i, i + \delta)$ does not
depend on $i$ for $\delta \in [4, t - 4]$, and the first, with $\delta = 4$ and $\delta = t - 4$, extends this to
$\delta = 3$ and $\delta = t - 3$ (here $t \ge 9$ is used). So $f(i, i + \delta) = f(\delta)$ with
$f(\delta - 1) = f(\delta + 1)$ for $\delta \in [4, t - 4]$, and $f(\delta) = f(t - \delta)$ because the two offsets
describe the same pair. As $t$ is odd, $\delta$ and $t - \delta$ have different parity, so $f$ is constant, say
$\alpha$, and $\rho_i = 2\alpha$.

\step{The recurrence.} Put $s_i = x(A_iB_i)$, $p_i = x(A_iB_{i+1})$, $q_i = x(A_iB_{i+2})$, $m_i = x(A_iB_{i-1})$,
$w_i = x(A_iB_{i-2})$, $b_i = x(B_iB_{i+1})$, $c_i = x(B_iB_{i+2})$, $e_i = x(A_iA_{i+2})$. The tight triples give
$p_i = 2\alpha - s_{i+1}$; $q_i = e_i + s_{i+2} = \alpha + m_{i+3} = p_i + b_{i+1}$; $x(A_iB_j) = \alpha + s_j$ and
$x(B_iB_j) = s_i + s_j + \alpha$ for $|i - j| \ge 3$; $c_i = w_{i+2} + s_{i+2} = m_{i+1} + p_{i+1} = b_i + b_{i+1}$;
$s_i + s_{i+3} + \alpha = w_{i+2} + p_{i+2}$. Hence $m_{i+3} = q_i - \alpha$, $b_{i+1} = q_i - 2\alpha + s_{i+1}$ and
$w_{i+2} = s_i + 2s_{i+3} - \alpha$, and the two expressions of $c_i = b_i + b_{i+1}$ through $m, p$ and through
$w, s$ become
\[
  q_{i-1} + q_i - q_{i-2} = 5\alpha - s_i - s_{i+1} - s_{i+2}, \qquad
  q_{i-1} + q_i = 3\alpha - s_{i+1} + s_{i+2} + 2s_{i+3},
\]
whence $q_i = -2\alpha + s_{i+2} + 2s_{i+4} + 2s_{i+5}$ and then
\begin{equation}\label{eq:rec}\tag{REC}
  s_i + 2 s_{i+3} + s_{i+4} = \tfrac72\,\alpha \qquad\text{for all } i \in \Z_t.
\end{equation}

\step{Rigidity.} The characteristic polynomial of \eqref{eq:rec} is $z^4 + 2z^3 + 1 = (z + 1)(z^3 + z^2 - z + 1)$.
The cubic has no rational root, so it is irreducible, and no cyclotomic polynomial has degree 3; hence its roots are
not roots of unity. For odd $t$, $-1$ is not a $t$-th root of unity either. So the circulant system \eqref{eq:rec}
on $\Z_t$ has only the constant solution $s_i = (7/8)\alpha$, and then $p, q, m, w, b, c, e$ and all far pairs are
determined by $\alpha$. The solution space of the tight system is a line, i.e.\ $C_t$ is extreme.
\end{proof}

\begin{proof}[Proof of (b)]
\step{Solutions on the line.} Let $\alpha \in \R$ and let $(s_j)$, $j \in \Z$, satisfy \eqref{eq:rec} for all
$i \in \Z$. Define $x$ on the pairs of $C_\infty$ by the formulas above ($x = \alpha$ on $A$-pairs at distance
$\ge 3$, $2\alpha$ at distance 1, and so on; the full list is in \file{family\_C\_proof.py}). Then $x$ satisfies
every tight triple of $C_\infty$: by the pattern argument it suffices to check the triples inside a window of 13
positions, and for each of these 397 triples the equation is a rational combination of instances of \eqref{eq:rec}
(check V2, exact). Let $\lambda \approx -1.8393$ be the real root of $z^3 + z^2 - z + 1$. Then $\alpha = 8$,
$s_j = 7 + \lambda^j$ is such a sequence; write $D(P) = (x(P) - d(P))/d(P)$ for the resulting relative deviation of a
pair $P$ of $C_\infty$. Call two pairs of $C_\infty$ \emph{related} if both are $A$-pairs at distance 1 or $\ge 3$,
or if they are $A_iB_j$ and $A_{i'}B_j$ with $|i - j|, |i' - j| \ge 3$.

\begin{lemmaS}
$D(P) = D(Q)$ only if $P = Q$ or $P$ and $Q$ are related.
\end{lemmaS}

\noindent\emph{Proof of Lemma S.} $D$ vanishes on the $A$-pairs at distance 1 or $\ge 3$, and
$D(A_iB_j) = \lambda^j/15$ for $|i - j| \ge 3$. For each of the other eight types of pairs, $D$ of the pair at
position $i$ is $\lambda^i\cdot\kappa$ with an explicit $\kappa \in \Q(\lambda)$, and all nine values of $\kappa$
(including $1/15$) are nonzero. An equality between two such pairs reads $\kappa = \lambda^k \kappa'$; since
$|\lambda| > 1$ this determines $k$ up to finitely many candidates, bounded rigorously by rational enclosures of
$\lambda$, and none of them gives equality in $\Q(\lambda)$ (check V3, exact). For a $B$-pair at distance $g \ge 3$,
$D = \lambda^i(1 + \lambda^g)/22$; an equality with a pair of another type reads
$\lambda^h(1 + \lambda^{-g}) = 22\kappa'$ for an integer $h$, where
$|1 + \lambda^{-g}| \in [1 - |\lambda|^{-3}, 1 + |\lambda|^{-3}]$ bounds $h$ and each $h$ leaves at most one $g$;
again no equality occurs (V3). For two different $B$-pairs at distances $\ge 3$,
$\lambda^i + \lambda^j = \lambda^{i'} + \lambda^{j'}$: if the largest exponent occurs on both sides the pairs
coincide, and otherwise the largest term exceeds the sum of the absolute values of the other three, because
$|\lambda|^{-1} + |\lambda|^{-3} + |\lambda|^{-4} < 0.8$. Finally $D \ne 0$ for every pair that is not an $A$-pair at
distance 1 or $\ge 3$. \hfill$\square$

\step{The partition.} Let $\pi$ be the partition of the pairs of $C_t$ whose classes are: the set of all $A$-pairs
at cyclic distance 1 or $\ge 3$; for each $j$ the set of pairs $A_iB_j$ with $|i - j| \ge 3$; and all other pairs as
singletons. $\pi$ is not total. We show that $\pi$ is closed under the merges of level $s$ when $t \ge 6s$; since
$\cl_s$ is the finest partition closed under these merges (Proposition~\ref{prop:closure}(i)), $\cl_s$ refines $\pi$
and is not total.

Let $W$ be a set of at most $s$ points of $C_t$. Its points occupy at most $s$ positions of $\Z_t$, which cut the
cycle into at most $s$ runs of free positions with $t - s \ge 5s$ positions in total, so some run of 5 consecutive
positions is free. Cutting the cycle inside this run identifies the positions of $W$ with integers in an interval
of $t - 5$ consecutive integers, and this identification $\psi\colon W \to C_\infty$ preserves distances: two points
at line offset $o$, $1 \le o \le t - 6$, have cyclic offset $o$ if $o < t/2$, and $o - t$ otherwise; in the second
case $o \ge 5$ and $t - o \ge 6$, so the pair is far (offset of absolute value at least 3) both on the line and on
the cycle, where the distance does not depend on the offset. Hence the tight triples inside $W$ are tight triples of
$C_\infty$, and for every real $c$ the function $u_c(P) = 1 + c\cdot D(\psi P)$ on the pairs inside $W$ is a
solution of $T_W$ in the normalised variables (a combination of the solutions $d$ and $x$). It is constant on the
classes of $\pi$: $D = 0$ on the first class, and $D(\psi(A_iB_j)) = \lambda^{\psi(j)}/15$ on the classes of the
second kind. So $u_c$ solves $T_W(\pi)$. If $W$ forced two different classes $C \ni P$, $C' \ni Q$ of $\pi$ to be
equal, we would have $u_c(P) = u_c(Q)$ for all $c$, i.e.\ $D(\psi P) = D(\psi Q)$, and by Lemma S the pairs
$\psi P$, $\psi Q$ would be equal or related, i.e.\ $P$ and $Q$ would lie in the same class of $\pi$. So no merge is
forced, $\pi$ is closed at level $s$, and $L(C_t) > s$ for every $s \le t/6$.
\end{proof}

\begin{proof}[Proof of (c)]
For $t \le 21$ the level-$(2t - 1)$ closure was computed in rational arithmetic and is not total
(\file{family\_C.py}). Let $t \ge 23$. As in (b) it suffices to show that $\pi$ is closed under the merges of level
$2t - 1$, i.e.\ that $W = V - v$ forces no two classes of $\pi$ to be equal, and by the rotation we may take
$v = A_0$ or $v = B_0$. For each of the two we construct a solution $x$ of $T_W(\pi)$ with $x = 0$ on the first
class of $\pi$ (a deviation from $d$, as $D$ in (b)) whose relative values $D(C) = x(C)/d(C)$ on the classes $C$ of
$\pi$ inside $W$ are pairwise different; then $u = 1 + cD$ separates any two classes. Indices are taken modulo $t$
and written with representatives near $0$.

\step{Step 1: parametrisation.} Let numbers $\sigma_j, q_j$ ($j \in \Z_t$) be given, and in addition $p_{-1}, p_0$
if $v = A_0$, and $w_{-1}$ if $v = B_0$. Put $p_i = -\sigma_{i+1}$ for the other $i$, and
\[
  m_i = q_{i-3}, \qquad b_i = q_{i-1} - p_{i-1}, \qquad w_i = \sigma_{i-2} + \sigma_{i+1} - p_i, \qquad
  c_i = b_i + b_{i+1}
\]
(for $v = B_0$: $w_{-1}$ as given, and $c_{-1} = w_1 + \sigma_1$). Define $x$ on the pairs inside $W$ by $x = 0$ on
the $A$-pairs at distance 1 or $\ge 3$, $x(A_iB_i) = \sigma_i$, $x(A_iB_{i+1}) = p_i$, $x(A_iB_{i+2}) = q_i$,
$x(A_iB_{i-1}) = m_i$, $x(A_iB_{i-2}) = w_i$, $x(B_iB_{i+1}) = b_i$, $x(B_iB_{i+2}) = c_i$,
$x(A_iA_{i+2}) = q_i - \sigma_{i+2}$, and $x(A_iB_j) = \sigma_j$, $x(B_iB_j) = \sigma_i + \sigma_j$ for
$|i - j| \ge 3$. (For $v = A_0$ the numbers $\sigma_0, q_0, p_0$ belong to pairs through $A_0$; they enter as
$\sigma_0 = x(A_iB_0)$ for $i$ far from $0$, $q_0 = x(A_3B_2)$ and $q_0 - p_0 = x(B_1B_2)$. For $v = B_0$ quantities
with an index that would give a pair through $B_0$ are not used, and $x(A_{-2}A_0)$ is arbitrary.) Consider the
equations
\[
  (5a_i)\quad c_i = w_{i+2} + \sigma_{i+2}, \qquad\qquad (5b_i)\quad c_i = m_{i+1} + p_{i+1}.
\]

\begin{lemmaP}
$x$ is constant on the classes of $\pi$, and it satisfies all tight triples inside $W$ if $5a_i$ and $5b_i$ hold for
all $i$ except $5a_{-2}$, $5b_{-1}$ when $v = A_0$, and except $5a_0$, $5a_{-1}$, $5a_{-2}$, $5b_0$, $5b_{-2}$ when
$v = B_0$.
\end{lemmaP}

\noindent\emph{Proof of Lemma P.} Go through the 11 types of tight triples in the proof of (a). The triples
$A_iA_{i+1}$ through $A_k$ read $0 = 0$; those through $B_{i+1}$ read $p_i + \sigma_{i+1} = 0$, and the two with
$i \in \{-1, 0\}$ (for $v = A_0$), resp.\ the one with $i = -1$ (for $v = B_0$), contain $v$. The three triples for
$A_iB_{i+2}$ are the definitions of $x(A_iA_{i+2})$, of $m_{i+3}$ and of $b_{i+1}$. $A_iB_j$ through $A_j$ reads
$\sigma_j = \sigma_j$. The triples for $B_iB_{i+2}$ are the definition of $c_i$ (through $B_{i+1}$), $5a_i$ (through
$A_{i+2}$) and $5b_i$ (through $A_{i+1}$). $B_iB_j$ through $A_i$ reads
$\sigma_i + \sigma_j = \sigma_i + \sigma_j$, and $B_iB_{i+3}$ through $A_{i+2}$ is the definition of $w_{i+2}$. The
triples that contain $v$ are exactly those in which the statement uses a free number or omits an equation: for
$v = A_0$, the triples $(A_{-1}, B_0, A_0)$, $(A_0, B_1, A_1)$, $(A_0, A_3, B_2)$, $(A_0, B_1, B_2)$,
$(A_i, A_0, B_0)$ free $p_{-1}$, $p_0$, $q_0$ and $\sigma_0$, and $(B_{-2}, A_0, B_0)$, $(B_{-1}, A_0, B_1)$ are
$5a_{-2}$, $5b_{-1}$; for $v = B_0$, the triple $(B_{-1}, B_0, B_1)$ frees $c_{-1}$, which is then defined by
$5a_{-1}$; $(B_{-3}, A_{-1}, B_0)$ frees $w_{-1}$; $(A_{-2}, A_0, B_0)$ is the only triple with the pair
$A_{-2}A_0$; and $5a_0$, $5a_{-2}$, $5b_0$, $5b_{-2}$ concern pairs through $B_0$. \hfill$\square$

\step{Step 2: the solution.} If $q_{i-2}, q_{i-1}, q_i$ are given by $q_j = \sigma_{j+2} + 2\sigma_{j+4} +
2\sigma_{j+5}$, the $p$'s by $p_i = -\sigma_{i+1}$, and $\sigma_{i+1} + 2\sigma_{i+4} + \sigma_{i+5} = 0$, then
$5a_i$ and $5b_i$ hold (the computation in the proof of (a)). Let $j_0 = 3$ for $v = A_0$ and $j_0 = 2$ for
$v = B_0$, $\eps = \lambda^{-t}$, and put
\[
  \sigma_j = g\cdot(-1)^j + a\cdot\lambda^{j-t} \quad\text{for } j_0 \le j \le t - 1, \qquad
  q_j = \sigma_{j+2} + 2\sigma_{j+4} + 2\sigma_{j+5} \quad\text{for } 8 \le j \le t - 8,
\]
with $(g, a) = (1 - \eps, 2)$ for $v = A_0$ and $(g, a) = (0, 1)$ for $v = B_0$. Both $(-1)^j$ and $\lambda^j$
satisfy the recurrence, so $5a_i$ and $5b_i$ hold for $11 \le i \le t - 9$. The remaining unknowns are $\sigma_j$ for
$0 \le j < j_0$ ($j = 1$ only, for $v = B_0$), $q_j$ for $|j| \le 7$, and $p_{-1}, p_0$ resp.\ $w_{-1}$: 20 resp.\
16 unknowns. The remaining equations are those $5a_i$, $5b_i$ with $-8 \le i \le 10$ that are required by Lemma P:
36 resp.\ 33 linear equations with rational coefficients, whose right-hand sides are of the form $r_0 + r_1\eps$
with $r_0, r_1 \in \Q(\lambda)$ that do not depend on $t$. They are solvable with solutions of the same form (check
W1 in \file{family\_C\_top.py}; for $v = A_0$ the system has rank 19 and one solvability condition, which is
$(1 - \eps)a = 2g$ and explains the choice of $g$; for $v = B_0$ the solution is unique). By Lemma P this gives a
solution $x$ of $T_W(\pi)$ for every odd $t \ge 23$. As a guard against slips in Lemma P, the solution was
instantiated on the real cycle for $t = 19, 21, 23, 25, 27, 31$ and every tight triple of $C_t - v$ was verified
exactly (W2). The same check confirms, class by class, that the symbolic values used in Step 3 are the values on the
cycle.

\step{Step 3: the values are pairwise different.} Away from the break the values of $x$ follow the formulas of (a)
with $s = \sigma$, so for a pair $P$ whose formula involves only $\sigma_j$ on the stretch, at position $t - k$
($k \ge 1$),
\[
  D(P) = \pm g\cdot\kappa' + a\cdot\kappa\cdot\lambda^{-k},
\]
where $\kappa \in \Q(\lambda)$ and $\kappa' \in \Q$ depend only on the type of the pair ($\kappa$ is the constant of
Lemma S, $\kappa'$ the corresponding constant for the root $-1$) and the sign on the parity of $k$. The pairs for
which this fails are 40 ($v = A_0$) resp.\ 27 ($v = B_0$) pairs near the break, with $D = r_0 + r_1\eps$, and the
far $B$-pairs, whose value is the sum of the values of $\sigma$ at their two positions divided by 22. So an equality
$D(P) = D(Q)$ between two classes is an equation
\[
  G_0 = \sum_i c_i\cdot\lambda^{-k_i}
\]
in at most five integer unknowns $k_i$ (the positions involved, and $t$ itself for the terms in $\eps$), with
$G_0, c_i \in \Q(\lambda)$ taken from a finite list: 2,305 equations for $v = A_0$ and 752 for $v = B_0$. Since
$|\lambda| > 1$ the solutions of such an equation can be listed completely: if $|G_0|$ exceeds the sum of the
largest possible values of the terms there is none; otherwise one branches on the largest term, fixing its exponent
at its least admissible value or raising that value; if $G_0 = 0$, one divides by the largest term. All bounds use
rational enclosures of $\lambda$, and equalities are decided exactly in $\Q(\lambda)$. The search (check W3) finds
109 resp.\ 53 solutions when the ranges of the $k_i$ are relaxed to $k_i \ge 1$, $t \ge 23$; in 92 resp.\ 38 of them
$P$ and $Q$ are the same class, and in the other 17 resp.\ 15 one of the pairs sits at one of the positions near
the break where the formula does not apply (these pairs are among the 40 resp.\ 27, with their true values). A pair
just left of the break whose formula involves only $\sigma$'s to the right of it (for example $A_{-1}A_1$, with
$\sigma_3$ and $\sigma_4$) belongs to the first kind with $k = t + 1$; the ranges of $k$ are set accordingly. So $D$
takes different values on different classes, for every odd $t \ge 23$. (For $v = B_0$ the pair $A_{-2}A_0$ is
excluded; it occurs in no tight triple inside $W$, so nothing forces its class to equal another.) Hence $W$ forces
no merge, $\pi$ is closed at level $2t - 1$, $\cl_{2t-1}$ refines $\pi$, and $L(C_t) = 2t$.
\end{proof}

The search procedure was tested on 300 equations with planted solutions (all recovered), and by duplicating one
class of each kind under a different name (every duplicate is reported as a coincidence); see \file{--selftest}.

\paragraph{Second implementation.} Checks W1--W3 were repeated by a second program written from the text of this
proof (\file{family\_C\_top\_check.py}; it shares no code with the first). It differs in three ways. (i) It fixes
$R = 10$ and treats every pair with an index within $R$ of the removed point as a class of its own with an explicit
value $r_0 + r_1\eps$, computed pair by pair from the formulas of Step 1 (354 resp.\ 334 such classes); all other
pairs have $R < k < t - R$, where the formula $D = \pm g\kappa' + a\kappa\lambda^{-k}$ applies for every type, so no
list of exceptional positions is needed. (ii) The equation $G_0 = \sum c_i \lambda^{-k_i}$ is solved by another
method: some term of a solution has absolute value at least $|G_0|/(\text{number of terms})$; this term and its
exponent are enumerated, and one recurses. (iii) A solution is discarded only if $P$ and $Q$ are the same class, if
an exponent violates $t - k > R$ (or the corresponding bound for a far $B$-pair), or if the two points of a far
$B$-pair are closer than 3; every other solution would be reported. Of 86,799 resp.\ 66,096 equations, 373 resp.\
400 have solutions (69 resp.\ 36 the same class, 304 resp.\ 364 out of range), and none is reported. The cut
solutions of the two programs are equal. The second program builds the solution on the real cycles
$t = 23, 25, 27, 29, 31, 37, 45$ with arrays indexed modulo $t$, verifies every tight triple, and checks that the
classes of $\pi$ inside $W$ correspond one to one, with equal values, to the instances of its list of classes that
are in range at $t$. Its search was compared with exhaustive enumeration in a box on 150 random equations (261
solutions, none missed), and it reports duplicated classes and classes planted to coincide at one value of $t$.

\begin{remarks}
(1) Part (b) is implied by (c); we keep it because its proof needs only Lemma S, while (c) relies on a larger
computation.
(2) The level-$(2t - 1)$ closure of $C_t$ has exactly $(t^2 + 13t + 2)/2$ classes for every odd $t$ from 9 to 21,
the number of classes of $\pi$: in these cases $\cl_{2t-1} = \pi$.
(3) Replacing $d(A_i, A_{i+2}) = 12$ by $16$ gives a second family, an extreme ray with $L = 2t$ for every odd $t$
from 9 to 21 by the same computation.
(4) How the family was found: the two-orbit metrics whose distances depend only on offsets capped at $\pm 3$ are,
for all large $t$, the points of a fixed 12-dimensional cone of ``patterns of radius 2'' (\file{pattern\_hunt.py});
of its 26,319 extreme rays with positive distances, 9,464 give extreme rays of $\MET_{2t}$ for at least two $t$
between 9 and 14, and 96 pass two filters (classes left unmerged on the cycle, and no merge forced on an open
stretch); 8 of the 96 have $L = 2t$ at $t = 9, 11, 13$, and they are $C_t$, its variant (3), their mirror images
and the versions with $A$ and $B$ exchanged.
(5) Chirality matters: the mirror-symmetric family $K_t$ of C6 reduces in the same way to a recurrence, but all its
characteristic roots are roots of unity, every solution on an open stretch is periodic, and its locality stays
bounded.
\end{remarks}

\section{A counting bound for small-choice generators}

\begin{avis}[Avis {\cite[Thm 3.3]{A80}}]
Let $0 < \gamma < 1/5$ and $n^{-1/5+\gamma} \le p \le 1 - n^{-1/4+\gamma}$. Then the shortest-path metric of
the random graph $G_{n,p}$ spans an extreme ray of $\MET_n$ with probability $1 - o(1)$.
\end{avis}

With $p = 1/2$ all labelled graphs are equally likely, so $2^{\binom n2}\cdot(1 - o(1))$ labelled graphs $G$ have
$d_G$ extreme \cite[p.~138]{A80}; this lower bound on the number of extreme rays of $\MET_n$ is quoted in this form
in \cite[\S5.1]{CL25}. The case $p = 1/2$ is also proved by Theorem~\ref{thm:four}(b), so what follows does not
depend on \cite[Thm 3.3]{A80}.

\begin{proposition}\label{prop:count}
The number $v(n)$ of orbits of vertices of $m_n$ satisfies
\[
  v(n) \ \ge\ \frac{2^{\binom n2}\cdot(1 - o(1))}{n!\cdot 2^{n-1}} \ =\ 2^{n^2/2 - O(n \log n)}.
\]
\end{proposition}

\begin{proof}
Distinct labelled graphs have distinct metrics $d_G$, and $G$ is recovered from $x_G = (2/\tau)d_G$ as the set of
pairs of smallest value, so the vertices $x_G$ of Lemma~\ref{lem:vertex} are distinct. An orbit has at most
$n!\cdot 2^{n-1}$ elements.
\end{proof}

\paragraph{Upper bound.} Cizma and Linial \cite[Thm 5.5]{CL25} prove that the number of faces of $\MET_n$ is at
most $2^{2.38\,n^2(1+o(1))}$. Hence the number of extreme rays of $\MET_n$, and the number of orbits of vertices of
$m_n$ adjacent to a cut, is $2^{\Theta(n^2)}$, with the constant in the exponent between $1/2$ and $2.38$.

\begin{definition}[generator]
A \emph{generator} consists of sets $O_n$ of \emph{objects on $n$ nodes} ($n \ge n_0$), a finite set $E$ of
exceptional objects, and $r$ \emph{rules}. Each rule maps a pair (object on $n - 1$ nodes, choice
$c \in \{1, \dots, C(n)\}$) to an object on $n$ nodes or to nothing. The generated objects are the exceptional ones
and the outputs of rules applied to generated objects. A map from generated objects on $n$ nodes to orbits of $m_n$
(possibly partial) says which orbit, if any, an object \emph{produces}. The \emph{choice} of the generator is
$b(n) = \log_2 C(n)$ bits.
\end{definition}

Nothing is assumed about the objects: they may be vertices, points of $m_n$ that are not vertices, graphs, or
partial structures. Restriction-keeping moves, families with boundedly many integer parameters of size polynomial
in $n$, and the generator of C9 for a fixed denominator are generators with $b(n) = O(\log n)$.

\begin{theorem}\label{thm:generator}
Let a generator have $b(n) \le (1 - \delta)n$ for some $\delta > 0$ and all large $n$. Then it produces at most
$2^{(1-\delta)n^2/2 + O(n)}$ orbits of $m_n$, and the share of the orbits of $m_n$ it produces tends to $0$. In
particular the number of orbits it misses is unbounded, so no finite list of exceptions completes it.
\end{theorem}

\begin{proof}
Let $N(n)$ be the number of generated objects on $n$ nodes. A generated object on $n$ nodes is exceptional or the
output of one of $r$ rules on one of $N(n - 1)$ objects with one of $2^{b(n)}$ choices, so
$N(n) \le r\cdot 2^{b(n)}\cdot N(n - 1) + |E|$, and by induction
\[
  N(n) \ \le\ (|E| + 1)\cdot\prod_{k \le n} \bigl(r\cdot 2^{b(k)} + 1\bigr) \ =\ 2^{\sum_{k \le n} b(k) + O(n)}.
\]
Since $b(k) \le (1 - \delta)k$ for all large $k$, $\sum_{k \le n} b(k) \le (1 - \delta)n^2/2 + O(n)$. The number of
orbits produced is at most $N(n)$. By Proposition~\ref{prop:count}, $v(n) \ge 2^{n^2/2 - O(n \log n)}$, so
$N(n)/v(n) \le 2^{-\delta n^2/2 + O(n \log n)} \to 0$ and $v(n) - N(n) \to \infty$. The exceptions are included in
$N(n)$.
\end{proof}

\begin{remarks}
(a) The bound is sharp in the choice up to a constant factor: with $b(n) = n - 1$ a rule can choose the
neighbourhood of the new node, hence an arbitrary graph, and then the single rule $G \mapsto x_G$ produces
$2^{n^2/2 - O(n \log n)}$ orbits; and by \cite{CL25} a choice of $4.76\cdot n$ bits per node is enough to index all
cut-adjacent orbits.
(b) The theorem does not need the $m_9$ list or any computation.
(c) It is a statement about growth, and the lists show the growth has begun:
$v(7), v(8), v(9) = 13,\ 533,\ \ge 1{,}056{,}367$, successive factors 41 and 1,982.
\end{remarks}

\begin{question}\label{q:one}
Let $r(n)$ be the number of extreme rays of $\MET_n$. Is $\log_2 r(n) \sim n^2/2$, i.e.\ are almost all extreme rays
graphic on the logarithmic scale, or is $\lim \log_2 r(n)/n^2 > 1/2$? By \cite{CL25} the limit, if it exists, is at
most $2.38$. On the lists the graphic share of the cut-adjacent orbits is 92\%, 20\%, 0.52\% for $n = 7, 8, 9$
(C10), which suggests that non-graphic rays dominate at least in number. (For vertices of $m_n$ not adjacent to a
cut we know only the trivial bound $2^{O(n^2 \log n)}$.)
\end{question}

\section{Computations: locality on the vertex lists and beyond}

All computations use exact rational arithmetic unless marked *. Lists: the vertices of $m_n$ for $n \le 8$
(complete: \cite{DFPS01,DFMV03}; completeness for $n = 8$ was re-proved independently in this project) and the list
of \cite{DI07} for $n = 9$ as distributed by its authors. The distributed file has 1,056,367 orbits including the
cut orbit; \cite{DI07} and the authors' web page state 1,056,368. The file has 1,056,369 lines of which two are
header lines, so the published number equals the number of lines minus one, which suggests a counting slip, but we
cannot exclude that one orbit is missing from the file. We verified that the 1,056,367 rows are vertices, pairwise
inequivalent, and include the 1,221 orbits of \cite{DI07} that are not adjacent to a cut. The $m_9$ list is not
known to be complete, so C2, C3 and C7--C12 are statements about the list.

\comp{C1 (all extreme rays of $\MET_k$, $k \le 8$).} Every extreme ray of $\MET_k$ with positive distances is
obtained from a vertex of $m_k$ by a switching (Lemma~\ref{lem:vertex} and its converse). Running through all
$2^{k-1}$ switchings of every orbit representative gives 1, 3, 33, 3,751 rays up to permutation for
$k = 5, 6, 7, 8$. As a check of completeness we counted all extreme rays of $\MET_n$ from these lists (a ray on $k$
points with automorphism group $A$ gives $k!\cdot S(n,k)/|A|$ rays on $n$ points by merging points,
\cite[Thm 4.2]{A80}; the cuts are the case $k = 2$): the totals 25, 296, 55,226, 119,269,588 for $n = 5, 6, 7, 8$
agree with the published values \cite{DDweb}, and the numbers of orbits under $S_n$ are 3, 8, 46, 3,918.
\begin{center}
\begin{tabular}{rrrrrrr}
\toprule
$k$ & rays & $L = 4$ & 5 & 6 & 7 & max $L$ \\
\midrule
5 & 1 & 1 & & & & 4 \\
6 & 3 & 3 & & & & 4 \\
7 & 33 & 32 & 1 & & & 5 \\
8 & 3,751 & 2,691 & 992 & 67 & 1 & 7 \\
\bottomrule
\end{tabular}
\end{center}
\emph{Validation of the implementation.} On 14,646 weighted graphs with 5 to 11 nodes (957 rays, 13,689 non-rays by
the rank test) the closure at $s = k$ ends with one class exactly for the rays, as
Proposition~\ref{prop:closure}(iii) requires.

\comp{C2 (the $m_9$ list).} For each of the 1,055,145 cut-adjacent orbits, one switching to a ray (the one with the
smallest alphabet); 3,394 of these rays have fewer than 9 distinct points and are counted with their own $k$.
\begin{center}
\begin{tabular}{rrrrrrrr}
\toprule
$k$ & rays & $L = 4$ & 5 & 6 & 7 & 8 & 9 \\
\midrule
5--7 & 81 & 81 & & & & & \\
8 & 3,313 & 2,411 & 870 & 27 & 5 & & \\
9 & 1,051,751 & 315,636 & 515,637 & 201,727 & 18,323 & 426 & 2 \\
\bottomrule
\end{tabular}
\end{center}
For $k = 9$ the shares are 30.0\%, 49.0\%, 19.2\%, 1.7\%, 0.04\% and two orbits with $L = 9$ (rows 1,048,423 and
1,048,435 of the list, both of denominator 27, weights 5 to 10, with 37 tight triples where 35 are needed). For
these two the closure stops changing at level 6 with 13 and 12 classes joined by 14 and 13 relations, one more than
the number of classes: the quotient system is a single circuit, and it has no two-term consequence on fewer than
all 9 points.

\comp{C3 (dependence on the switching).} $L$ is an invariant of the ray, not of the orbit of $m_n$: different
switchings of a vertex give different rays. For $n \le 8$ all switchings are covered by C1. For 4,000 random orbits
of the $m_9$ list and all 256 switchings (27,826 rays, 27,737 of them on 9 distinct points) the distribution over
the latter is $L = 4$: 30.4\%, 5: 48.4\%, 6: 19.7\%, 7: 1.5\%, 8: 0.03\%, close to the distribution of C2. Per
orbit, the smallest $L$ over the switchings is 4 for 49\% of the orbits, 5 for 42\%, 6 for 9\%, 7 for 0.2\%; the
largest is 4 for 5\%, 5 for 43\%, 6 for 46\%, 7 for 7\%, 8 for 0.2\%; the two differ by at most 4. So 9\% of the
sampled orbits have no switching with $L \le 5$.

\comp{Reading of C1--C3.} The largest locality is 4, 4, 5, 7, 9 for $k = 5, \dots, 9$, and the share of rays with
$L = 4$ decreases (97\%, 72\%, 30\% for $k = 7, 8, 9$). Fixed-scope deduction characterises the extreme rays for
$k \le 6$ (scope 4) and $k \le 7$ (scope 5) only, and at $k = 9$ there are rays whose extremality has no proof from
any proper subset of the points. Large locality is rare on the lists: two orbits in a million at $k = 9$.

\comp{C4 (random and circulant graphs).} Random graphs $G(n, 1/2)$, 300 samples for each $n$: see the remark after
Theorem~\ref{thm:four}. Every extreme ray found among the weighted circulant graphs on 7 to 13 nodes with weights at
most 3 has $L = 4$. High locality is not typical, and it is not a consequence of symmetry alone.

\comp{C5 (searches for large locality beyond the lists).} Values marked * were computed with a symmetry-reduced
closure modulo $2^{31} - 1$ and not re-run in rational arithmetic.

(a) \emph{Two-orbit metrics.} The metrics on points $A_i, B_i$ ($i \in \Z_t$) invariant under the rotation form a
cone whose extreme rays can be enumerated and tested. Cyclic symmetry: $t = 5$ (294 rays of $\MET_{10}$), largest
$L = 8$; $t = 6$ (3,019 rays of $\MET_{12}$), largest $L = 9$; $t = 5$ with a fixed centre (1,332 rays of
$\MET_{11}$), largest $L = 8$. Dihedral symmetry: $t = 7$ (3,497 rays of $\MET_{14}$), largest $L = 12$; $t = 7$
with a fixed centre ($n = 15$), largest $L = 13$; $t = 9$ (55,300 rays of $\MET_{18}$ tested, of 132,051
candidates), largest $L = 17$*.

(b) \emph{Generic procedures do not find large locality.} Seven vertices of $m_{10}$ found by random walks on the
skeleton (49 rays over all switchings): $L$ from 4 to 7. Adding a point to the rays of $\MET_8$ with $L \ge 6$ (beam
search to 11 points, 6,794 rays): $L$ never above 7. Random face descent in $\MET_n$, $n = 7, \dots, 11$ (54 extreme
rays from 280 descents, the others ended with a zero distance): $L \le 5$. The graphs $K_1 \vee C_m^2$ (a centre
joined to the square of the $m$-cycle) are graphic extreme rays with $L = 7$ for every $m$ from 8 to 13.

\comp{C6 (a symmetric family whose locality saturates).} The family $K_t$ on $A_i, B_i$ ($i \in \Z_t$):
$d(A_i, A_{i\pm1}) = 10$, $d(A_i, A_{i\pm2}) = 20$, $d(A_i, A_j) = 26$ otherwise; $d(B_i, B_{i\pm1}) = 24$,
$d(B_i, B_j) = 12$ otherwise; $d(A_i, B_i) = 7$, $d(A_i, B_{i\pm1}) = 17$, $d(A_i, B_{i\pm2}) = 13$,
$d(A_i, B_j) = 19$ otherwise. It is invariant under the dihedral group. By computation $K_t$ is an extreme ray for
$t = 7, 11, 13, 17, 19$ and for no other $t$ from 7 to 20, with $L = 12$ for $t = 7$, $L = 16$* for $t = 11$ and
$14 < L \le 16$* for $t = 13$. After the 4-point merges the classes are one class $K$ for all pairs of $B$'s and
singletons, and eliminating all but $s_i = [A_iB_i]$ leaves the recurrence
$s_i + s_{i+1} + 2s_{i+2} + s_{i+3} + s_{i+4} = 6K$ for all $i$ mod $t$. Its symbol
$(1 + \omega^2)(1 + \omega + \omega^2)$ vanishes exactly at the primitive 3rd and 4th roots of unity. So $K_t$
cannot be an extreme ray when 3 or 4 divides $t$, and every solution along an open stretch of the cycle has period
12, so 13 consecutive positions force $s_j = s_{j+12}$; this caps the locality of the family. (The failure for
$t = 10, 14$ is not explained by this recurrence alone.) Of 1,094 mirror-symmetric local-pattern families obtained
by extending the dihedral rays of C5(a) for $t = 7$, every one has a window of 14 positions that forces merges
beyond the 4-point ones. This is what led to the chiral families of Section~\ref{sec:main}.

\comp{Largest locality known per $n$.} 4, 4, 5, 7 for $n = 5, \dots, 8$ (complete lists); 9 for $n = 9$ (the $m_9$
list); at least 8, 8, 9, 12, 13 for $n = 10, 11, 12, 14, 15$ (C5); $n$ for every $n \equiv 2 \pmod 4$, $n \ge 18$
(Theorem~\ref{thm:main}(c)).

\begin{fact}
$L(d) \le n - 1$ for every extreme ray of $\MET_n$, $n \le 8$, but $\max L = n$ for $n = 9$.
\end{fact}

\begin{question}
Is $\max\,\{ L(d) : d \text{ extreme in } \MET_n \} = n$ for every $n \ge 9$? By Theorem~\ref{thm:main}(c) this
holds for $n \equiv 2 \pmod 4$, $n \ge 18$, and for $n = 9$ by C2 (the two rays exist whether or not the $m_9$ list
is complete).
\end{question}

\begin{question}
Does the share of orbits of $m_n$ with $L = 4$ tend to $0$? By Theorem~\ref{thm:four} and
Proposition~\ref{prop:count} at least $2^{n^2/2 - O(n \log n)}$ orbits have $L = 4$, so this is
Question~\ref{q:one} in another form.
\end{question}

\section{Other descriptions of the $m_9$ list}

\comp{C7 (primitive vertices; cut adjacency).} Call a vertex of $m_n$ \emph{primitive} if none of its restrictions
to $k$ nodes, $3 \le k < n$, is a vertex of $m_k$. A rule that keeps its input as a restriction of its output cannot
produce a primitive vertex.
\begin{center}
\begin{tabular}{rrrr}
\toprule
$n$ & orbits & primitive & share \\
\midrule
7 & 13 & 4 & 31\% \\
8 & 533 & 219 & 41\% \\
9 & 1,056,367 & 673,787 & 63.8\% \\
\bottomrule
\end{tabular}
\end{center}
On the $m_9$ list 1,221 orbits are not adjacent to a cut, in agreement with \cite{DI07}.

\comp{C8 (a normal form).} Say node $v$ can be \emph{peeled} from the node set $S$ of a vertex $x$ if $x|_S$ is a
vertex of the fiber of $m_S$ over $x|_{S - v}$, i.e.\ the coordinates $x_{va}$ are determined by the tight
inequalities inside $S$ that contain $v$. Peel until no node can be peeled; the remaining set is a \emph{seed}. On
the $m_9$ list the seed has 3 nodes for 583,636 orbits, 4 nodes for 472,722 and 5 nodes for 8. A tight triangle
through $v$ gives an equation $x_{va} \pm x_{vb} = c$; these form a signed graph on $S - v$, and $v$ can be peeled
iff every component has an unbalanced cycle. Describing a step by the lengths of the shortest unbalanced cycles of
its components, the 1,056,366 non-cut orbits use six kinds of steps: one 3-cycle (5,760,177 steps), two components
with 3-cycles (88,230), one 5-cycle (11,967), one 2-cycle (3,497), one 7-cycle (1,565), a 3-cycle and a 5-cycle
(22).

\comp{C9 (the normal form is not a construction).} The generator ``start from all points of $m_3$ and $m_4$ with
denominator $D$, extend by all determined steps with denominator $D$, test for vertices'' finds every orbit of
$m_8$ with $D = 3, 4, 5, 6$ (52, 12, 15, 95 orbits) and nothing else, from 2,958, 67,164, 757,052 and 5,215,846
candidates on 8 nodes, of which 43\%, 0.28\%, 0.04\%, 0.07\% are vertices. The partial points that lie on a path to
a vertex number 39, 92, 200, 287 on 4 to 7 nodes for $D = 6$ (against 85, 311, 1,871, 22,788 generated). Six local
features of an extension were tested as filters that never discard a useful extension; the best keeps 86--100\% of
the extensions where 17--56\% are useful.

\comp{C10 (graph-indexed families).} Alphabets and underlying graphs of the cut-adjacent orbits:
\begin{center}
\begin{tabular}{rrrrrr}
\toprule
$n$ & orbits & alphabets & graphic & underlying graphs & most orbits on one graph \\
\midrule
7 & 12 & 2 & 92\% & 11 & 1 \\
8 & 532 & 9 & 20\% & 256 & 12 \\
9 & 1,055,145 & 1,273 & 0.52\% & 31,061 & 2,177 \\
\bottomrule
\end{tabular}
\end{center}
A graph does not determine its weights, so ``graph $\mapsto$ vertex'' needs the weights as input.

\comp{C11 (the simplest weighted family).} For rays with alphabet $\subseteq \{1, 2\}$ (1, 3, 28, 925 for
$k = 5, \dots, 8$; 27,963 from the $m_9$ list) a deduction rule in the spirit of \cite[\S2]{A80} (opposite sides of
geodesic rectangles; odd closed chains of exchanges; propagation through triples with one unknown) is sufficient for
extremality and is necessary for all 957 rays with $k \le 8$; it fails for 76 of the 27,963 rays on 9 nodes. The
best necessary condition found rejects 92\% of 7,344 near-miss non-rays. This repeats, for weights $\{1, 2\}$, the
experience of \cite[pp.~132--133]{A80}.

\comp{C12 (elliptope).} For $x$ in an orbit let $\lambda(x)$ be the smallest eigenvalue of $J - 2X$, an orbit
invariant. On $m_7$, $m_8$ and 1,000 sampled orbits of $m_9$, $\lambda$ does not separate the orbits reached by 8
restriction-keeping rules from the others (probability that an unreached orbit has smaller $\lambda$ than a reached
one: 0.43, 0.47; 0.46 at equal denominators).

\section{Related work}

Avis \cite{A80} introduced graphic extreme rays, the isometric-cycle criterion, the density theorem used here and
the embedding theorem (every rational metric is a restriction of an extreme ray, \cite[Thm 3.5]{A80}); the last
explains why ``is a restriction of a graphic vertex'' cannot characterise vertices. Laurent and Poljak \cite{LP92}
and Laurent \cite{L96} study graphic and one-third-integral vertices and the operations that preserve vertices. The
lists are due to Deza, Deza, Fukuda, Pasechnik, Sato, Mizutani, Vo and Indik \cite{DDF96,DFPS01,DFMV03,DI07}. Cizma
and Linial \cite{CL25} count path systems and faces of $\MET_n$: they quote Avis's lower bound
$2^{\binom n2(1+o(1))}$ for the extreme rays, prove the upper bound $2^{2.38n^2(1+o(1))}$ for all faces, and show
that metrizability of a path system is decided by a linear program; as far as we could see they do not consider
generators or bounded-scope criteria. Proposition~\ref{prop:count} is thus Avis's bound transferred to orbits of
$m_n$, and Theorem~\ref{thm:generator} is an elementary consequence of it. What is new here, to our knowledge, is
the locality invariant with Theorems~\ref{thm:four} and~\ref{thm:main}, the formulation of the counting bound as a
limit on generators, and the computations.

Bandelt and Dress \cite{BD92} decompose a metric into split metrics and a split-prime residue; extreme rays other
than cuts are split-prime, but split-prime metrics need not be extreme, so this gives a necessary condition only.
Koolen, Moulton and T\"onges \cite{KMT00} and Sturmfels and Yu \cite{SY04} classify the six-point metrics. Lavrov
\cite{La25} generalises Avis's isometric-cycle criterion to metrics on simplicial posets and states that the
classification of extreme metrics for general $n$ is open. Bodwin \cite{Bo18} characterises the path systems that
are unique shortest paths of a weighted graph by forbidden subsystems coming from two-coloured polyhedra; his
examples have 4 to 10 nodes and the family is not bounded in size. That result concerns faces of the metric cone of
a different kind (strict metrizability), but it is the closest analogue we know of Theorem~\ref{thm:main}:
obstructions that cannot be confined to a bounded number of points. Melleray, Petrov and Vershik \cite{MPV08} study
extreme rays under the name of linearly rigid metric spaces; they pose as an open question to describe the extremal
metric spaces on $n$ points and to estimate their number. Casabella, Joswig and Kastner \cite{CJK24} classify the
metric spaces on seven points through subdivisions of hypersimplices and give new infinite families of rays.
Feldman and Kehoe \cite{FK26} give further families of extreme rays (bowtie metrics). In none of these did we find
the locality invariant, a bound on generators, or a statement that extreme rays admit no bounded-scope
characterisation.

\emph{Forward citations.} We went through the works citing \cite{A80} (68), \cite{L96} (26) and \cite{DI07} (6) in
the OpenAlex index (November 2026) and read those whose titles concern extreme rays, vertices of the metric polytope
or criteria for extremality; none contains a result of the kind of Theorems~\ref{thm:main}
or~\ref{thm:generator}.

\section{Reproduction}

Scripts in \file{scripts/phase9/}, outputs in \file{results/phase9/}.

{\small\raggedright
\begin{center}
\begin{tabular}{@{}>{\raggedright\arraybackslash}p{0.21\textwidth}>{\raggedright\arraybackslash}p{0.37\textwidth}>{\raggedright\arraybackslash}p{0.36\textwidth}@{}}
\toprule
item & script & output \\
\midrule
Theorem~\ref{thm:main}(a), (b): checks V1--V3; remark 4 &
  \file{family\_C\_proof.py}, \file{family\_C\_structure.py}, \file{pattern\_hunt.py}, \file{cand\_pattern.py} &
  \file{family\_C\_proof.json}, \file{pattern\_hunt.json}, \file{cand\_pattern.json} \\
Theorem~\ref{thm:main}(c): $t \le 21$; checks W1--W3 for $t \ge 23$ &
  \file{family\_C.py}; \file{family\_C\_top.py} & \file{family\_C.json}; \file{family\_C\_top.json} \\
Theorem~\ref{thm:main}(c): second implementation of W1--W3 &
  \file{family\_C\_top\_check.py} & \file{family\_C\_top\_check.json} \\
C1, validation & \file{locality\_checks.py}, \file{ray\_count\_check.py} &
  \file{locality\_checks.json}, \file{ray\_count\_check.json} \\
C2 & \file{locality.py --all9} & \file{locality\_m9all.json} \\
C3 & \file{locality\_switch9.py} & \file{locality\_switch9.json} \\
C4 & \file{random\_graph\_locality.py}, \file{locality\_checks.py} &
  \file{random\_graph\_locality.json}, \file{locality\_checks.json} \\
C5 & \file{sym\_search.py} (\file{--mirror 0 1 --fast} for the dihedral runs), \file{locality\_checks.py},
  \file{locality\_climb.py}, \file{random\_ray\_locality.py} &
  \file{sym\_search.json}, \file{locality\_checks.json}, \file{locality\_climb.json},
  \file{random\_ray\_locality.json} \\
C6 & \file{famK.py}, \file{famK\_level.py}, \file{family\_scan.py}, \file{window\_test.py}, \file{exact\_check.py} &
  \file{family\_scan.json}, \file{window\_test.json}, \file{exact\_check.json} \\
C7, C8 & \file{vertex\_structure.py}, \file{step\_types.py} &
  \file{vertex\_structure\_deza9all.json}, \file{step\_types\_m9all.json} \\
C9 & \file{fiber\_generator.py}, \file{useful\_share.py}, \file{choice\_features.py} &
  \file{fiber\_generator\_n8.json}, \file{useful\_share\_n8.json}, \file{choice\_features\_n8.json} \\
C10 & \file{graph\_families.py --all9} & \file{graph\_families\_m9all.json} \\
C11 & \file{family12.py}, \file{judge\_family.py} & \file{round\_18\_family.json} \\
C12 & \file{elliptope.py} & \file{elliptope.json} \\
\bottomrule
\end{tabular}
\end{center}
}

Ranks are computed modulo $2^{31} - 1$ with an exact fallback whenever the modular rank is deficient, so every ``is
a vertex'' and ``is a ray'' verdict is exact; the locality closure uses rational arithmetic throughout, except for
the values marked * in C5 and C6 and the parenthesis of Theorem~\ref{thm:main}(c).

\begin{thebibliography}{DFMV03}

\bibitem[A80]{A80} D. Avis, On the extreme rays of the metric cone, \emph{Canad. J. Math.} 32 (1980) 126--144.

\bibitem[BD92]{BD92} H.-J. Bandelt, A. W. M. Dress, A canonical decomposition theory for metrics on a finite set,
\emph{Adv. Math.} 92 (1992) 47--105.

\bibitem[Bo18]{Bo18} G. Bodwin, On the structure of unique shortest paths in graphs, arXiv:1804.09745.

\bibitem[CJK24]{CJK24} L. Casabella, M. Joswig, L. Kastner, Subdivisions of hypersimplices: with a view toward
finite metric spaces, \emph{Exp. Math.} (2024), arXiv:2402.17665.

\bibitem[CL25]{CL25} D. Cizma, N. Linial, On the number of path systems, arXiv:2508.21379v3 (2025).

\bibitem[DD-web]{DDweb} A. Deza, metric polytope data page, \url{www.cas.mcmaster.ca/~deza/metric.html}.

\bibitem[DDF96]{DDF96} A. Deza, M. Deza, K. Fukuda, On skeletons, diameters and volumes of metric polyhedra,
\emph{LNCS} 1120 (1996) 112--128.

\bibitem[DFMV03]{DFMV03} A. Deza, K. Fukuda, T. Mizutani, C. Vo, On the face lattice of the metric polytope,
\emph{LNCS} 2866 (2003) 118--128.

\bibitem[DFPS01]{DFPS01} A. Deza, K. Fukuda, D. Pasechnik, M. Sato, On the skeleton of the metric polytope,
\emph{LNCS} 2098 (2001) 125--136.

\bibitem[DI07]{DI07} A. Deza, G. Indik, A counterexample to the dominating set conjecture, \emph{Optim. Lett.} 1
(2007) 163--169.

\bibitem[DL97]{DL97} M. Deza, M. Laurent, \emph{Geometry of Cuts and Metrics}, Springer, 1997.

\bibitem[FK26]{FK26} D. V. Feldman, E. R. Kehoe, Extreme rays of the metric cone and the metric body: bowtie
metrics and half-one extremality, preprint, Zenodo, doi:10.5281/zenodo.21898516 (2026).

\bibitem[KMT00]{KMT00} J. Koolen, V. Moulton, U. T\"onges, A classification of the six-point prime metrics,
\emph{European J. Combin.} 21 (2000) 815--829.

\bibitem[La25]{La25} A. Lavrov, Ternary relations and their polytopes, arXiv:2509.06811 (2025).

\bibitem[L96]{L96} M. Laurent, Graphic vertices of the metric polytope, \emph{Discrete Math.} 151 (1996) 131--153.

\bibitem[LP92]{LP92} M. Laurent, S. Poljak, One-third-integrality in the max-cut problem, \emph{Math. Programming}
71 (1995) 29--50 (CWI Report BS-R9209, 1992).

\bibitem[MPV08]{MPV08} J. Melleray, F. V. Petrov, A. M. Vershik, Linearly rigid metric spaces and the embedding
problem, \emph{Fund. Math.} 199 (2008) 177--194, arXiv:math/0611049.

\bibitem[SY04]{SY04} B. Sturmfels, J. Yu, Classification of six-point metrics, \emph{Electron. J. Combin.} 11
(2004) R44.

\end{thebibliography}

\end{document}
